% MATH Companion Part VI / VI-30 / explanatory bridge only.
\paragraph{From ring dimension to scheme dimension.}
For a scheme \(X\), dimension is defined using strict chains of
irreducible closed subsets:
\[
 Z_0\subsetneq Z_1\subsetneq\cdots\subsetneq Z_r.
\]
One does \emph{not} need to embed \(X\) into an affine space to define
this quantity.  When \(X=\operatorname{Spec}A\), the correspondence
between primes and irreducible closed sets shows directly that
\(\dim X=\dim A\).  For a general scheme, affine charts allow the
same notion to be studied locally, while the topology records how
the charts are joined.

\paragraph{What changes when we pass to a local ring.}
A point \(x\in X\) with affine neighbourhood
\(U=\operatorname{Spec}A\) corresponds to a prime \(\mathfrak p\).
Its stalk is \(\mathcal O_{X,x}=A_{\mathfrak p}\), and thus
\[
 \dim\mathcal O_{X,x}=\operatorname{ht}(\mathfrak p).
\]
This measures how long a chain of specializations can end at \(x\).
It need not equal \(\dim X\) for an arbitrary point: at the
generic point \(\eta\) of the affine line,
\(\mathcal O_{\mathbb A^1_k,\eta}=k(t)\) is a field of
dimension zero, even though \(\dim\mathbb A^1_k=1\).

\paragraph{Why closed points behave differently on a variety.}
Suppose now that \(X\) is an \emph{integral scheme of finite type}
over a field \(k\).  Let \(x\) be a closed point, and take an
affine neighbourhood \(\operatorname{Spec}A\) of \(x\).
The corresponding prime \(\mathfrak m\) is maximal; by Zariski's
lemma \(A/\mathfrak m\) is a finite extension of \(k\).
Hence \(\dim(A/\mathfrak m)=0\).  The dimension formula for a
finite-type integral \(k\)-algebra now yields
\[
 \underbrace{\operatorname{ht}(\mathfrak m)}_{\dim\mathcal O_{X,x}}
 +\underbrace{\dim(A/\mathfrak m)}_{0}
 =\dim A=\operatorname{trdeg}_k K(X).
\]
Nonempty affine opens in an integral finite-type variety have the
same function field and dimension.  Therefore
\(\dim\mathcal O_{X,x}=\dim X\).  The finite-type and
integrality hypotheses are part of this reasoning, not optional
decoration.

\paragraph{A calculation with two descriptions of the same curve.}
Consider the cuspidal curve
\[
 C=\operatorname{Spec}k[t^2,t^3]
 \simeq \operatorname{Spec}\bigl(k[x,y]/(y^2-x^3)\bigr).
\]
Its ring is a one-dimensional finitely generated domain; its
function field is \(k(t)\).  Consequently,
\[
 \dim C=\operatorname{trdeg}_k k(t)=1.
\]
At the cusp, the maximal ideal \((t^2,t^3)\) has height one, so
the local ring also has dimension one.  Nevertheless the tangent
space at the cusp has dimension two (calculated in VI/32).
The comparison teaches an important lesson: \emph{local Krull
dimension} and \emph{tangent-space dimension} are different
invariants and can disagree at singular points.

\paragraph{Why nonempty open subsets retain dimension in the integral case.}
If \(U\subseteq X\) is a nonempty open subset of an integral
finite-type \(k\)-variety, then \(U\) contains a nonempty affine
open \(V\).  The rings of nonempty affine opens in \(X\) have the
same function field \(K(X)\), and each has dimension
\(\operatorname{trdeg}_k K(X)\).  This proves that \(\dim U=\dim X\).
The argument cannot simply be transplanted to disconnected or
reducible schemes.  For example,
\[
 X=\operatorname{Spec}k\;\amalg\;\mathbb A^1_k
\]
has dimension one, but the nonempty open-and-closed component
\(\operatorname{Spec}k\) has dimension zero.  Thus hypotheses
must remain visible wherever a global dimension identity is used.
